\emph{Setting.} $\Cl(1,3)$ with $\gamma_{0}^{2} = 1$, $\gamma_{i}^{2} = -1$, $I_{4} = \gamma_{0}\gamma_{1}\gamma_{2}\gamma_{3}$, $I_{4}^{2} = -1$, $\tilde I_{4} = I_{4}$. For a versor $V$ with $V\tilde V$ a nonzero scalar, $V^{-1} = \tilde V/(V\tilde V)$ and $\mathrm{Ad}(V)x = VxV^{-1}$ preserves $Q$. The wedge is $W = \{x^{1} > |x^{0}|\}$, $B = \gamma_{0}\gamma_{1}$, $D = I_{4}B$.

\emph{(i).} $B^{2} = \gamma_{0}\gamma_{1}\gamma_{0}\gamma_{1} = -\gamma_{0}^{2}\gamma_{1}^{2} = 1$ and $\tilde B = -B$, so $B\tilde B = -1$ and $B^{-1} = B$: $B$ is a versor of the full spin group outside the identity component. $B$ anticommutes with $\gamma_{0}$ and $\gamma_{1}$ and commutes with $\gamma_{2}$, $\gamma_{3}$, so $\mathrm{Ad}(B)$ negates $x^{0}$, $x^{1}$ and fixes $x^{2}$, $x^{3}$: it is $j_{W}$, and it maps $W$ to $\{x^{1} < -|x^{0}|\}$, the causal complement $W'$. Next $D = I_{4}\gamma_{0}\gamma_{1} = \gamma_{2}\gamma_{3}$ (moving $\gamma_{0}\gamma_{1}$ through), $D^{2} = -1$, so $\mathrm{Ad}(D)$ is the turn by $\pi$ in the plane $\gamma_{2}\gamma_{3}$, and $\mathrm{Ad}(I_{4})x = -x$ for every vector since $I_{4}$ anticommutes with vectors. As $I_{4}D = I_{4}^{2}B = -B$, $\mathrm{Ad}(B) = \mathrm{Ad}(I_{4}D) = \mathrm{Ad}(I_{4}) \circ \mathrm{Ad}(D)$. That the turn by $\pi$ in the dual plane is the boost continued over half a modular period, differing from it by $x \mapsto -x$ on vectors, is \claimref{synthesis.modular-period-on-values-is-the-dual-plane-turn}(ii), cited, not re-proved. The Bisognano-Wichmann conjugation of a Wightman field is CPT composed with the turn by $\pi$ about the boost axis, acting on spacetime as $j_{W}$ (Fact \claimref{angle.imported-theorems-deliver-their-hypotheses}); (i) gives its geometric part as PT after the dual-plane turn. The script checks every identity and the action on three vectors.

\emph{(ii).} The geometric part is proved here and the modular action is imported. $Bn_{\pm} = \mp n_{\pm}$ with $n_{\pm} = \pm\gamma_{0} + \gamma_{1}$, so $\mathrm{Ad}(B)n_{\pm} = -n_{\pm}$; $\gamma_{2}$, $\gamma_{3}$ are fixed. For the boost rotor $R = \cosh(\eta/2) + B\sinh(\eta/2)$, $R n_{\pm}\tilde R = R^{2}n_{\pm} = e^{\mp\eta}n_{\pm}$, checked exactly at $R = (5 + 4B)/3$ ($R\tilde R = 1$, $e^{\eta} = 9$), and $R$ maps a point of $W$ into $W$. The translations by $a n_{+}$ and $a n_{-}$, $a \ge 0$, map $W$ into itself ($x^{1} + a > |x^{0} \pm a|$ when $x^{1} > |x^{0}|$). Under the spectrum condition (the generators of the future null translations are positive) and if the vacuum is cyclic and separating for the wedge algebra $M$, Borchers' theorem \cite{Borchers1992} (Fact \claimref{synthesis.imported-kuckert-modular-pct-spin-statistics-and-borchers-theorem}(ii); its modular-group form also in Fact \claimref{angle.imported-theorems-deliver-their-hypotheses}) gives $\Delta^{it}U(a)\Delta^{-it} = U(e^{-2\pi t}a)$ and $JU(a)J = U(-a)$ on each half-sided translation group; the transverse translations map $M$ onto itself and fix the vacuum, so they commute with $\Delta$ and $J$. Hence the modular objects act on all translations as the boost at rapidity $-2\pi t$ and as $j_{W}$, which is the geometric action of (i) and (ii); nothing in the argument uses free fields. The spectrum condition of the interacting net is a hypothesis here (Kuckert's Assumption I); the store's (N$_{+}$) is posed for quasi-free states of a CAR mode algebra on horizon generators (\claimref{synthesis.positive-null-energy-gives-r-premises}) and does not supply it. The conjugation half of Borchers' theorem is cited from memory.

\emph{(iii).} Let $R_{\theta} = a + b\gamma_{1}\gamma_{2}$ with $a^{2} + b^{2} = n$; $R_{\theta}\tilde R_{\theta} = n$ and $\mathrm{Ad}(R_{\theta})$ turns $\gamma_{1}$ toward $\gamma_{2}$ by $\theta$ with $\cos\theta = (a^{2} - b^{2})/n$, $\sin\theta = 2ab/n$. Since $\gamma_{0}$ commutes with $\gamma_{1}\gamma_{2}$, $B_{\theta} = R_{\theta}B\tilde R_{\theta}/n = \gamma_{0}(\cos\theta\,\gamma_{1} + \sin\theta\,\gamma_{2})$, the boost bivector of $W_{\theta} = \mathrm{Ad}(R_{\theta})W$, and $B_{\theta}B = -\gamma_{0}^{2}(\cos\theta\,\gamma_{1} + \sin\theta\,\gamma_{2})\gamma_{1} = \cos\theta + \sin\theta\,\gamma_{1}\gamma_{2} = R_{\theta}^{2}/n$, the spin rotor whose $\mathrm{Ad}$ is the turn by $2\theta$ (checked at five Pythagorean half-angles with the double-angle action on $\gamma_{1}$). At $\theta = \pi$, $R_{\pi} = \gamma_{1}\gamma_{2}$, $B_{\pi} = -B$, $W_{\pi} = W'$ and $B_{\pi}B = -1$; at $\theta = 0$ the product is $1$. The path $\theta \mapsto \cos\theta + \sin\theta\,\gamma_{1}\gamma_{2}$, $\theta \in [0, \pi]$, is the lift starting at $1$ of the loop of turns by $2\theta$, a closed loop of rotations, and it ends at $-1$: the rotation by $2\pi$ of the spin group, $-1$ on spinor values and $(-1)^{2s}$ on spin-$s$ values. As maps $\mathrm{Ad}(B_{\pi}) = \mathrm{Ad}(-B) = \mathrm{Ad}(B)$, so the sign at $\theta = \pi$ is fixed only by the continuity of the lift. A form that does not depend on any choice of sign: $B_{\pi/2} = \gamma_{0}\gamma_{2}$ and $(\epsilon_{1}B_{\pi/2}\,\epsilon_{2}B)^{2} = (\gamma_{0}\gamma_{2}\gamma_{0}\gamma_{1})^{2} = (\gamma_{2}\gamma_{1})^{2} = -1$ for $\epsilon_{1}, \epsilon_{2} = \pm 1$, so the reflections of $W$ and of its quarter-turned copy compose to a rotor of the half turn whose square is the rotation by $2\pi$ in the spin group (checked for all four sign choices, with $\mathrm{Ad}$ of the product negating $\gamma_{1}$). In Kuckert's theorem the sign that matters is carried at the field level by the twist (iv). In Kuckert's notation $R(\varphi) = J_{\varphi/2}J_{W}$ \cite{Kuckert1995} (Corollary 3.2; Fact \claimref{synthesis.imported-kuckert-modular-pct-spin-statistics-and-borchers-theorem}(i)), the product of the conjugations of the turned wedge and of $W$; (iii) is its geometric content: the product of the reflections is the lift of the rotation by $\varphi$, continuous in $\varphi$ and $-1$ at $\varphi = 2\pi$.

\emph{(iv).} This is the standard Klein-twist identity of the field net with $i$ written as $I_{4}$; the script represents the grading by even involutions of $\Cl(1,3)$ commuting with $I_{4}$ and the conjugate-linear $J$ by the reversal of $I_{4}$. $I_{4}$ is central in the even subalgebra, so a real span of $1$, $I_{4}$, $k$, $I_{4}k$ is commutative. With $k^{2} = 1$: $(1 + I_{4}k)^{2} = 1 + 2I_{4}k + I_{4}^{2}k^{2} = 2I_{4}k$ and $(1 + I_{4})^{2} = 2I_{4}$, so $V^{2} = k$; $(1 + I_{4})^{-1} = (1 - I_{4})/2$. The adjoint reverses the unit and fixes $k$: $V^{*} = (1 - I_{4}k)(1 - I_{4})^{-1}$, and $VV^{*} = (1 - I_{4}^{2}k^{2})/(1 - I_{4}^{2}) = 1$. A conjugate-linear $J$ reverses $I_{4}$ again and fixes $k$, so $JV^{*}J = V$ and $VJV^{*}J = V^{2} = k$. The script checks these for three even involutions $k$ commuting with $I_{4}$ and that $V \ne k$. This is the step $R(2\pi) = J_{\pi}J_{\mathcal{F}} = VJ_{\mathcal{F}}V^{*}J_{\mathcal{F}} = V^{2} = k$ of \cite{Kuckert1995} (Corollary 3.5), where $J_{\pi} = VJ_{\mathcal{F}}V^{*}$ comes from twisted duality, which Kuckert derives (from \cite{DoplicherRoberts1990}) rather than assumes, and $VJ_{\mathcal{F}} = J_{\mathcal{F}}V^{*}$; with the field theory's unit $I_{4}$ in place of $i$.

\emph{(v).} Imported: under Kuckert's Assumptions I to III (covering-group covariance with the spectrum condition, a unique cyclic vacuum, irreducibility, separability; modular P1CT symmetry $J_{W}\mathcal{A}(O)J_{W} = \mathcal{A}(j_{W}O)$ for double cones; a compact group of internal symmetries) in dimension $1 + s$, $s \ge 2$, every covariant normal field has $R(2\pi) = k$ \cite{Kuckert1995}; the paper notes that the modular groups need not act geometrically, so the premise is weaker than modular covariance, which implies it (Bisognano and Wichmann for Wightman fields \cite{BW1976}; Guido and Longo, Fact \claimref{synthesis.imported-free-field-spin-statistics-and-guido-longo-theorem}(iii)). For the series' free fields the same sign is already fixed by locality with positive norm (\claimref{synthesis.locality-of-the-attached-free-fields-fixes-their-statistics}); the theorem matters for interacting fields.

\emph{(vi).} Argued. The premise has a translation part and an algebra part. The translation part follows by (ii) from the spectrum condition, which is a hypothesis of the theorem and is not supplied by the store for an interacting net: (N$_{+}$) concerns quasi-free states of a CAR mode algebra, and the positive conformal energy result concerns the representation on a given state's GNS space. The algebra part, $J_{W}\mathcal{A}(O)J_{W} = \mathcal{A}(j_{W}O)$ for all double cones, is not implied: Borchers' relations constrain $J_{W}$ only through its commutation with the translations, and the KMS property at period $2\pi$ for the boosts, which (N$_{+}$) gives for quasi-free states, identifies the modular group on the wedge algebra but says nothing about where $J_{W}$ sends the algebras of regions inside $W$ (and \claimref{synthesis.modular-period-on-values-is-the-dual-plane-turn} shows that KMS at $2\pi$ does not fix statistics). Neither (HS), a statement about two-point functions on generators, nor (HD), about the state's occupation profile, bears on it. The field net of the theorem carries charges localisable in double cones or spacelike cones; the electric charge, constrained by the Gauss law, is not localisable there, so the electron's spin-statistics connection beyond perturbation theory needs more than this premise. Transport-absent: the theorem is standard; the reading of the reflection as the PT inversion after the half modular period and the twist with the unit $I_{4}$ re-express it, as the stored modular-period result classifies its own $I_{4}$ reading.
